\section{Benchmark and Audit Protocol}

\subsection{CRAFT-Core}

The benchmark begins with three executable domains:
\begin{itemize}[leftmargin=1.5em]
\item algebra,
\item graph path,
\item blocksworld-like planning.
\end{itemize}
We verbalize each problem into natural-language reasoning traces and construct
matched counterfactual families around the audited locus. The benchmark focuses
on executable, auditable settings because the goal is to isolate verifier
failure modes cleanly before making stronger natural-language generalization
claims.

\subsection{Counterfactual families}

The core families are:
\begin{itemize}[leftmargin=1.5em]
\item \texttt{local-invalid}: the process is changed locally near the audited
locus;
\item \texttt{answer-swap}: the process is kept fixed while the final answer
surface changes;
\item \texttt{lucky-answer}: an invalid process is paired with a correct-looking
final answer;
\item \texttt{paraphrase}: the language surface changes while semantics are
preserved.
\end{itemize}
\texttt{Delayed-repair} is treated as a stress slice rather than part of the
core benchmark, because it mixes local invalidity with later trajectory
recovery.

\begin{figure}[t]
\centering
\setlength{\fboxsep}{10pt}
\fbox{
\begin{minipage}{0.92\linewidth}
\small
\textbf{Executable latent problem} $\rightarrow$ \textbf{audited locus} $\rightarrow$
\textbf{matched quartet}

\vspace{0.5em}
\begin{tabular}{ll}
\texttt{valid\_correct} & locally valid process, correct answer surface \\
\texttt{invalid\_correct} & locally invalid process, same answer surface \\
\texttt{valid\_swapped} & same valid process, swapped answer surface \\
\texttt{invalid\_swapped} & same invalid process, swapped answer surface
\end{tabular}

\vspace{0.75em}
Audit intent:
\begin{itemize}[leftmargin=1.5em]
\item \amcd: compare valid vs.\ invalid while holding the answer surface fixed
\item \asstotal: compare the same process under answer-visible and answer-swapped views
\item artifact audit: reject shallow stylistic cues before using the benchmark for claims
\end{itemize}
\end{minipage}
}
\caption{Schematic of the matched-quartet audit protocol. The benchmark is built
from executable latent trajectories, then expanded into answer-matched local
counterfactuals so that process validity and answer consistency can be probed
separately.}
\label{fig:quartet_protocol}
\end{figure}

\subsection{Artifact audit}

Every benchmark stage is filtered through artifact checks before it is used for
claims. The initial hard slice already passes shallow-feature audit without
obvious style leakage. The cleaner benchmark is then accepted under a narrower
artifact-focused criterion: the goal is an artifact-clean audited benchmark,
not universally human-indistinguishable valid/invalid pairs. Under that
criterion, the current \textsc{Craft-Clean} export retains \texttt{52 / 54}
accepted families and \texttt{208} records. External human blind audit on this
export is still pending, so the benchmark should be read as strongly audited
rather than fully closed.
Throughout the paper, this artifact-focused criterion is the primary
benchmark-acceptance rule; the stricter universal detectability criterion is
reported only as a diagnostic.

\subsection{Two benchmark regimes}

We report two benchmark regimes on purpose rather than by accident:
\begin{itemize}[leftmargin=1.5em]
\item \textbf{\textsc{Craft-Deploy}} is the main deployment benchmark. It is a
naturalized full-hybrid generated slice designed to ask whether a verifier or
reranker improves selection quality and reduces exploitability under realistic
distribution shift.
\item \textbf{\textsc{Craft-Clean}} is the cleaner audit benchmark. It is the
accepted cleaner export under the artifact-focused acceptance criterion
and is more useful for isolating process-faithfulness gaps than for maximizing
utility separation.
\end{itemize}
These regimes should not be merged into a single headline score. They answer
rather different questions, and that difference is now part of the core claim
rather than a nuisance variable. Figure~\ref{fig:quartet_protocol} shows the
core logic shared by both regimes.

\subsection{External-source corroboration with \textsc{ProcessBench}}

The matched-quartet benchmarks are not our only evidence. To reduce dependence
on self-generated source problems, we also build an external-source benchmark
from \textsc{ProcessBench}. This line is deliberately not forced into the
quartet template. Instead, it provides two companion tasks:
\begin{itemize}[leftmargin=1.5em]
\item \textbf{PB-Trace:} whole-trace process validity prediction using the
annotated first-incorrect-step label,
\item \textbf{PB-Prefix:} prefix-level detection in which prefixes before the
first incorrect step are valid and prefixes from the first incorrect step
onward are invalid.
\end{itemize}
This external-source package is deliberately lighter than the matched-quartet
protocol, so it is not a substitute for the full \amcd/\ass{} benchmark.
Instead, it serves two corroboration roles. First, \textbf{PB-Trace} and
\textbf{PB-Prefix} test whether the main ranking story survives on a benchmark
built from externally sourced step-level supervision. Second, we add two small
external-source counterfactual checks directly on \textsc{ProcessBench}: an
answer-only swap set for external \ass{}-style stress, and an answer-matched
local-repair subset for minimal local-discrimination stress.
To keep this line reproducible, we materialize it as a canonical benchmark
package rather than a loose collection of scripts: the release now includes a
suite summary, a manifest of authoritative datasets and artifacts, and a split
manifest that pins the grouped \textbf{PB-Trace} split and the capped
\textbf{PB-Prefix} split used by the current baselines. In other words,
\textsc{ProcessBench} is now closed for its intended corroboration role even
though it is not a matched-quartet replacement.
As we show later, these corroboration lines support the same conservative
read. Visible answer sensitivity persists on real traces, but the answer-visible
method story still does not recover, and the strongest current verifier on
this external-source benchmark remains a simple step-only frozen baseline
rather than a visible reranker.
